% The same synthetic manuscript after acting on the findings.
%
% Every measured value in sample-manuscript.tex survives here. Nothing was
% deleted to make a finding go away: the recital advisories were answered by
% saying what the quantities establish, so that a passage is no longer an
% uninterrupted run of numeral-bearing sentences, and the cohesion advisories
% by carrying a noun forward rather than by adding connectives that fake a link.
%
% Numbers preserved: 4200, 0.013, 0.0004, 12, 10, 200, 8, 1.2, 4.6, 5, 0.00,
% 0.40, 350, 0.02, 4, 3.5, 0.7, 0.263, 4.2, 0.002, 0.0002, 0.041, 0.004,
% 0.028, 2, 0.007, 3, 0.052, 0.0011, 200, 0.0008, 0.0003, 0.0005.
\documentclass{article}

\newcommand{\Nsims}{4200}
\newcommand{\mbias}{0.013}
\newcommand{\cbias}{0.0004}

\begin{document}

\title{Multiplicative Shear Bias in Stacked Image Simulations}
\author{A. Sample}
\maketitle

\begin{abstract}
We measure the multiplicative shear bias of a moments-based shape estimator on
\Nsims{} image simulations. The stacked bias is $m = \mbias$ with an additive
term $c = \cbias$. That stacked bias is an average over a population, and it
hides where the estimator actually fails. Blending sets the residual bias, not
the signal-to-noise ratio, and at a blending fraction of 0.40 the bias reaches
0.052.
\end{abstract}

\section{Introduction}

Weak lensing surveys measure the coherent distortion of background galaxy
shapes. That distortion is small, so the estimator that measures it has to be
calibrated to a precision well below the survey's own statistical error.
Calibrating any estimator requires a population whose true shear is known,
which is what image simulations supply. This work calibrates a moments-based
estimator against such a population and reports where the calibration stops
holding.

The calibration is a chain, and the interesting failure is at its end. A
simulated population is built, its distribution is matched to the observed
catalog, the estimator is run over it, and a bias model is fitted to what comes
back. Fitting that model is routine. Propagating its residual into the survey
error budget is not, and that residual is what this paper is about.

\section{Method}

The simulation grid spans signal-to-noise ratio, size, and blending fraction,
since those are the three axes along which a moments estimator is known to
fail. Signal-to-noise ratio is sampled in 12 bins from 10 to 200, size in 8
bins from 1.2 to 4.6 pixels, and blending fraction in 5 steps from 0.00 to
0.40. Each cell of that grid holds 350 realizations at an input shear of 0.02
applied in 4 position angles, which is \Nsims{} images in total.

The images themselves are conventional. A Moffat point spread function with
$\beta = 3.5$ and a full width at half maximum of 0.7 arcsec is rendered on a
0.263 arcsec pixel scale, with 4.2 electrons of read noise.

Conventional images are measured with an unconventional estimator. It takes a
weighted second moment rather than fitting a profile, so it recovers the
quadrupole without assuming a light distribution, and its weight function is a
circular Gaussian matched to the object size. That size matching is the source
of the bias this paper calibrates, because a size-matched weight responds to
the object's size as well as to its shape.

\section{Results}

The stacked bias is $m = \mbias \pm 0.002$ with $c = \cbias \pm 0.0002$. That
stacked figure averages over a population in which the bias varies by an order
of magnitude. Sorted by signal-to-noise ratio that bias runs from $m = 0.041$ in
the lowest bin to $m = 0.004$ in the highest, and sorted by size the same bias
runs from 0.028 below 2 pixels to 0.007 above 3.

Neither of those trends limits the calibration. What limits the calibration is
blending, whose bias reaches 0.052 at a blending fraction of 0.40. Correcting
for that blending leaves a residual of 0.0011 in every cell of the grid.

That residual is a property of the estimator rather than of the sample, and
three resamplings say so. The three resamplings are a bootstrap over 200
resamples, which scatters by 0.0008, a split of the grid by position angle,
which moves $m$ by 0.0003, and a fresh noise realization, which moves it by
0.0005. All three are smaller than the residual they were meant to explain.

\section{Discussion and Conclusion}

The estimator is calibrated to $m = \mbias$, and blending rather than
simulation depth is what limits that calibration. Deeper simulations will not
remove a residual that survives every resampling of the present ones. Denser
ones will, if they sample the blending fraction beyond 0.40.

Two assumptions bound that conclusion. The photometric redshift distribution
was held fixed, and detector nonlinearity was not modelled. Neither assumption
enters the blending term, so neither changes the limit above; both set how far
the quoted $m$ transfers to a survey whose point spread function differs from
the one simulated here.

\end{document}
